\documentclass{article}
\usepackage{amsmath, amssymb, amsthm}
\usepackage{hyperref}
\usepackage{geometry}
\geometry{margin=1in}

\title{Erd\H{o}s Problem \#143}
\date{Last updated: 2025-08-31}
\author{Source: erdosproblems.com}

\begin{document}
\maketitle

\noindent\textbf{Status:} OPEN \quad \textbf{Prize: \$500}

\noindent\textbf{Tags:} primitive sets

\medskip
\noindent\textbf{Problem Statement:}

Let $A\subset (1,\infty)$ be a countably infinite set such that for all $x\neq y\in A$ and integers $k\geq 1$ we have\[ \lvert kx -y\rvert \geq 1.\]Does this imply that $A$ is sparse? In particular, does this imply that\[\sum_{x\in A}\frac{1}{x\log x}<\infty\]or\[\sum_{\substack{x <n\\ x\in A}}\frac{1}{x}=o(\log n)?\]

\bigskip
\noindent\textbf{Remarks:}

Note that if $A$ is a set of integers then the condition implies that $A$ is a primitive set (that is, no element of $A$ is divisible by any other), for which the convergence of $\sum_{n\in A}\frac{1}{n\log n}$ was proved by Erd\H{o}s \cite{Er35}, and the upper bound\[\sum_{\substack{x <n\\ x\in A}}\frac{1}{n}\ll \frac{\log x}{\sqrt{\log\log x}}\]was proved by Behrend \cite{Be35}. This $O(\cdot)$ bound was improved to a $o(\cdot)$ bound by Erd\H{o}s, S\'{a}rk\H{o}zy, and Szemer\'{e}di \cite{ESS67}.

In \cite{Er73} and \cite{Er77c} Erd\H{o}s mentions an unpublished proof of Haight that\[\lim \frac{\lvert A\cap [1,x]\rvert}{x}=0\]holds if the elements of $A$ are independent over $\mathbb{Q}$.

Over the years Erd\H{o}s asked for various different quantitative estimates, for example\[\liminf \frac{\lvert A\cap [1,x]\rvert}{x}=0\]or even (motivated by Behrend's bound)\[\sum_{\substack{x <n\\ x\in A}}\frac{1}{x}\ll \frac{\log x}{\sqrt{\log\log x}}.\]In \cite{Er97c} he offers \$500 for resolving the questions in the main problem statement above.

This was partially resolved by Koukoulopoulos, Lamzouri, and Lichtman \cite{KLL25}, who proved that we must have\[\sum_{\substack{x <n\\ x\in A}}\frac{1}{x}=o(\log n).\]See also [858].

\bigskip
\noindent\textbf{References:}

[Be35] Behrend, F., On sequences of numbers not divisible by another. London Math. Soc. Journal (1935), 42-45.
 
 [ESS67] Erd\H{o}s, P. and S\'{a}rk\"ozy, A. and Szemer\'{e}di, E., On a theorem of Behrend. J. Austral. Math. Soc. (1967), 9--16.
 
 [Er35] Erd\H{o}s, Paul, Note on Sequences of Integers No One of Which is Divisible By Any Other. J. London Math. Soc. (1935), 126-128.
 
 [Er73] Erd\H{o}s, P., Problems and results on combinatorial number theory. A survey of combinatorial theory (Proc. Internat. Sympos., Colorado State Univ., Fort Collins, Colo., 1971) (1973), 117-138.
 
 [Er77c] Erd\H{o}s, Paul, Problems and results on combinatorial number theory. III. Number theory day (Proc. Conf., Rockefeller Univ.,
New York, 1976) (1977), 43-72.
 
 [Er97c] Erd\H{o}s, Paul, Some of my favorite problems and results. The mathematics of Paul Erd\H{o}s, I (1997), 47-67.
 
 [KLL25] D. Koukoulopoulos, Y. Lamzouri, and J. D. Lichtman, Erd\H{o}s's integer dilation approximation problem and GCD graphs. arXiv:2502.09539 (2025).

\bigskip
\noindent\small{Source: \url{https://www.erdosproblems.com/143}}
\end{document}
